% Transcription par lecture visuelle directe, sans OCR.
% Statut : transcrit ; relecture indépendante à effectuer.
% Source : PDF original, vue 101, côté droit, page manuscrite 203.
\chapter{Suite de la préfecture de M. Philpin}
\begin{center}
Commencement de la persécution dirigée contre le pape.\\
Nouvelles épreuves pour la Congrégation, second bref d'indulgences.
\end{center}

\section*{Commencement de la persécution}
% Mention marginale des pages manuscrites 203–204.
Vous savez, Messieurs, que l'Église de la terre est appelée l'Église militante, il ne faut donc pas s'étonner si nous la voyons constamment occupée à soutenir de nouveaux combats. Si parfois Dieu veut bien lui accorder quelques moments de relâche, ce n'est qu'afin qu'elle puisse mettre à profit ces courts intervalles de repos pour se disposer à supporter avec plus de force et de courage les nouveaux assauts qui vont lui être livrés.

Il n'y avait encore que peu d'années que nous avions vu renaître la paix et déjà nous allons avoir à vous entretenir d'une affreuse tempête. L'Église sera battue de toute part par les vagues qui s'amoncèlent autour d'elle ; nous savons d'avance qu'elle triomphera de toute leur fureur, car ayant été solidement établie par J. C. lui-même, étant appuyée sur l'ancre inébranlable des promesses qu'elle a reçues de son divin fondateur et placée sous le pavillon glorieux de la croix avec le garant assuré de la victoire, elle vogue à travers les flots irrités d'une mer en courroux, toujours certaine de résister à tous les dangers qui la menacent ; mais la Congrégation, cette humble et fragile nacelle amarrée à la barque de Pierre par les chaînes bien solides de l'esprit de foi et de charité, que va-t-elle devenir ? Les promesses de perpétuité n'ont pas été faites pour elle. Ne devons-nous pas craindre de la voir disparaître dans l'épouvantable orage qui va fondre sur nous ?
% Source : PDF original, vue 102, côté gauche, page manuscrite 204.
Rassurez-vous, Messieurs, c'est sous la protection de Marie Vierge immaculée que cette société a été placée. Marie est là. Marie est notre mère. Elle saura bien défendre notre… J'allais dire sa Congrégation, car elle lui appartient. Sans doute il faut nous attendre à la voir en apparence abandonnée à toute la fureur des vents, cruellement secouée par la tempête. Un instant même elle paraîtra s'enfoncer sous les eaux de la tribulation ; on croira que c'en est fait d'elle… qu'elle est submergée pour toujours. Mais celui qui a conservé Jonas dans le ventre de la baleine, la conservera aussi dans cet état de mort apparente, et quand il en sera temps la rappellera à la vie. Attendez un peu ; vous la verrez bientôt après qu'elle aura avec l'Église souffert la persécution dans la personne de plusieurs de ses membres, se relever glorieuse et triomphante avec elle et reprendre à sa suite, avec plus de zèle que jamais, le cours accoutumé de ses bonnes œuvres.

Entrons donc dans le récit de ces événements ; ils se rattachent si directement à l'histoire de la Congrégation, qu'ils ne nous en écarteront pas. Nous y trouverons de nouvelles occasions de glorifier notre Dieu et d'observer encore les effets merveilleux de sa miséricorde qui ne cesse pas de se manifester sur nous même au milieu des châtiments que nos infidélités nous ont justement mérités.

Nous avons déjà raconté comment l'ambition avait placé Napoléon sur le trône ; il est dans la nature de cette passion de n'être jamais satisfaite. Napoléon éprouva les tristes effets de ce désir insatiable d'honneurs et de gloire. Il se jeta dans une voie fatale qui le conduisit à sa perte, bien loin de conserver
% Source : PDF original, vue 102, côté droit, page manuscrite 205.
quelques sentiments de reconnaissance envers le souverain pontife, il conçut le coupable désir de lui enlever toute puissance temporelle, et en même temps de subordonner à sa volonté suprême, l'exercice de la puissance spirituelle. De plus Napoléon avait pris la résolution de répudier sa femme ; il exigeait que cet acte scandaleux réprouvé par les lois de l'Église fût légitimé par le St Père. Sa politique mettait à cette sanction une haute importance, car il méditait de s'allier avec une de ces races antiques dans lesquelles le monde, depuis des siècles, avait continué à trouver ses maîtres. Il s'était persuadé qu'il serait impossible au pape, qu'il savait foible et impuissant, de résister un seul instant à la volonté d'un conquérant, chef d'une grande nation, et qui d'un clin d'œil pouvait mettre un million d'hommes en mouvement. Mais le St Père qui dans sa charité avait poussé la condescendance aussi loin qu'il l'avait cru possible, trouva dans le sentiment de son devoir la fermeté du martyre, dès que sa conscience lui fit une loi de la déployer.

\section*{Expulsion des Pères de la Foi}
% Mention marginale de la page manuscrite 206 ; reportée au début du récit.
L'empereur irrité de cette résistance, se prépara à user de tous les moyens que son despotisme orgueilleux mettait en sa puissance pour en triompher. Toutefois il redoutait encore l'opinion de la France. Il se rappelait comment le St Père y avait été reçu et malgré l'ivresse dont le remplissait la pensée de son immense pouvoir, il ne pouvait pas se dissimuler qu'il aurait à lutter contre les dispositions d'un peuple nombreux et d'un clergé fidèle, sincèrement attaché à la religion et qui s'efforcerait de soutenir la cause du chef de l'Église. Dans sa politique toute démoniaque, il voulut du moins, s'il ne pouvait le faire disparaître entièrement, affoiblir cet obstacle. Pour y parvenir
% Source : PDF original, vue 103, côté gauche, page manuscrite 206.
il commença à se mettre en garde contre ceux qui par leurs dispositions bien connues et la confiance dont ils étaient en possession, paraissaient pouvoir embrasser plus vivement le parti du pontife opprimé et être les plus capables de confirmer les fidèles dans leurs sentiments de soumission à l'autorité du vicaire de J. C. Parmi ceux-ci se montraient en première ligne les Pères de la Foi. Déjà ils avaient établi plusieurs collèges en France, deux de ces établissements étaient dans le diocèse de Lyon, l'un dans la petite ville de Belley, qui depuis est devenue la ville épiscopale du diocèse de ce nom ; l'autre était dans la ville de Roanne. L'archevêque leur avait en outre confié le petit séminaire de l'Argentières. Ces religieux successeurs des jésuites étaient comme eux inviolablement attachés au chef de l'Église. Par un quatrième vœu ajouté aux trois vœux ordinaires de religion, ils promettaient, comme leurs prédécesseurs, obéissance passive au souverain pontife. Napoléon ne pouvait donc pas douter que dans les projets hostiles qu'il méditait il les aurait pour adversaires ; leur opposition pouvait d'autant plus nuire à ses desseins, que soit au moyen de leurs maisons d'éducation, soit par la confiance générale qu'ils inspiraient, ils avaient les moyens d'agir fortement sur l'opinion publique. Napoléon pour enlever cette ressource à sa victime, lance un décret par lequel il supprime l'ordre des Pères de la Foi, et ordonne la clôture de tous leurs établissements. Cette mesure fut exécutée rigoureusement et avec rapidité comme l'étaient toujours les ordres du despote. Nous eûmes la douleur de voir fermer ces maisons qui commençaient à peine à donner à la religion
% Source : PDF original, vue 103, côté droit, page manuscrite 207.
l'espérance de la dédommager bientôt par la génération naissante, des maux qu'elle avait soufferts. Oh ! combien de parents chrétiens furent vivement affligés, en se voyant enlever la ressource qui leur offrait une sécurité entière pour l'éducation de leurs enfants. Déjà la maison de Belley qui avait été établie la première et dont l'existence s'est prolongée pendant quelques années avait rendu de très grands services. La Congrégation en particulier doit une grande reconnaissance aux saints religieux qui mettaient tous leurs soins à y former la jeunesse, car plusieurs de ses meilleurs sujets sont sortis de cet excellent établissement.

Nous déplorions donc avec amertume comme tous ceux qui étaient sincèrement attachés à la religion, ces pertes cruelles qui nous avaient de nouveau plongés dans l'état de crainte et d'inquiété dont nous étions à peine sortis. Quant aux vénérables pères, objets de cette persécution, à l'exemple des apôtres, ils se réjouissaient d'avoir été jugés dignes de souffrir quelque chose pour le nom de J. C. Ils opérèrent leur déménagement avec la même sérénité que s'ils se fussent disposés à une partie de plaisir et abandonnèrent leurs maisons en chantant le Te Deum en actions de grâces. Les mesures de rigueur ne s'arrêtèrent pas là. Nous avions espéré que le gouvernement satisfait d'avoir expulsé les Pères de la Foi de leurs collèges, leur laisserait du moins la liberté de choisir ainsi que tous les Français en ont le droit, le lieu qui leur conviendrait pour leur résidence. Quel bonheur c'eût été pour les nombreux fidèles qui avaient placé en eux leur confiance de pouvoir encore être encouragés par leurs conseils et édifiés par leurs exemples. Le clergé séculier aurait trouvé en eux d'utiles auxiliaires pour la prédication et l'administration
% Source : PDF original, vue 104, côté gauche, page manuscrite 208.
des sacrements. Vain espoir ! La politique ombrageuse de Napoléon craignait trop leur influence pour les laisser ainsi en repos au milieu d'une population dont ils étaient également aimés et estimés. L'ordre est donné à chacun d'eux à aller dans le diocèse auquel il appartient par sa naissance, se mettre à la disposition de son évêque.

\section*{Départ du père Roger}
% Mention marginale de la page manuscrite 208.
Voilà donc le P. Roger forcé de quitter Lyon pour se rendre à Coutances, sa patrie. Il faut qu'il abandonne toutes ces Congrégations dont il était chéri comme un père… Quel coup pour chacune d'elles ! Quel coup plus cruel encore peut-être pour la nôtre en particulier ! Quelle douleur malgré toute sa résignation, pour l'âme de ce bon père. Lui qui avait véritablement pour eux des entrailles paternelles, il faut qu'il délaisse ses enfants, ses enfants chéris, que lui avait lui-même rassemblés autour de lui dans un nouveau paradis terrestre, au milieu duquel il avait été chargé par le ciel d'implanter comme l'arbre de vie, cet esprit de foi qui seul peut vivifier nos actions, nos vœux, nos désirs et jusqu'à nos pensées ! Ah ! Messieurs, vous devez juger combien fut pénible pour les uns et pour les autres ce douloureux sacrifice. Encore n'en connaissions-nous pas alors toute l'étendue. Nous nous étions flattés que cette séparation ne serait qu'une épreuve passagère, Dieu sans doute le permettait ainsi ; car si nous avions pu dès lors prévoir la trop longue durée de cette absence, cette seule pensée eût suffi pour jeter le trouble et le désordre dans les cœurs des congréganistes, y exciter le découragement et peut-être ébranler la Congrégation. Le P. Roger lui-même espérait aussi revenir prochainement. Pour que nous ne souffrissions de son
% Source : PDF original, vue 104, côté droit, page manuscrite 209.
absence, il chargea l'un de ses confrères, M. Boissard qui avait été supérieur du collège de Roanne, de prendre soin de la Congrégation. Ce ne fut qu'après avoir pourvu autant qu'il était en lui à ce qui nous concernait qu'il partit pour Coutances. Il y fut placé à la tête d'un petit séminaire, et trouva aussi un nouvel aliment à son zèle pour la gloire de Dieu. Nous le verrons plus tard se fixer à Paris et y poursuivre le bien des congréganistes, en attendant le moment fortuné où la Providence réunira enfin le père et ses enfants.

\section*{M. Bochard, directeur de la Congrégation}
% Mention marginale de la page manuscrite 209.
Le P. Roger en nous mettant entre les mains du P. Boissard croyait que son confrère plus heureux que lui ne serait pas obligé de s'éloigner ; je ne me rappelle plus sur quel motif il avait fondé cet espoir ; ce qu'il y a de certain est qu'il ne se réalisa pas. Peu de temps après le départ du P. Roger, le P. Boissard fut également obligé de quitter la ville. Le cardinal crut alors devoir confier la direction de la Congrégation à M. Bochard, l'un de ses vicaires généraux. M. Bochard était un vertueux prêtre rempli de zèle et de piété ; il avait le plus grand désir de conserver le bien, mais il ne sut malheureusement comprendre ni l'esprit qui avait fondé notre Congrégation, ni celui qui la dirigeait. Il se présenta à nous plutôt en maître qu'en père. Son caractère était froid et son esprit s'arrêtait facilement devant les minuties. En général, autant le P. Roger s'était efforcé d'élargir les cœurs et les esprits, autant il nous semble que pour la direction qu'il cherchait à nous imprimer, M. Bochard aurait eu de la tendance à les rétrécir. Tout cela ressemblait si peu au P. Roger que ce n'était pas de nature à faire naître la confiance. Aussi le nouveau directeur fut-il reçu très froidement par la Congrégation. On eut pour lui des égards et de la déférence comme on les lui devait en sa qualité de
% Source : PDF original, vue 105, côté gauche, page manuscrite 210.
délégué de l'autorité ecclésiastique, mais les sympathies des congréganistes vis-à-vis de lui ne purent jamais aller plus loin que cette soumission purement passive ; ils n'en persévérèrent pas moins dans un état de défiance continuelle.

\section*{Mort de M. Deroche}
% Mention marginale des pages manuscrites 210–212.
Pendant que la Congrégation gémissait ainsi sur la perte douloureuse de son père, et sur la position délicate où elle était par suite de ce grand sacrifice, Dieu qui voulait qu'elle fût alors éprouvée de toutes les manières, lui envoya encore une nouvelle affliction. Pour la première fois, la mort vint éclaircir nos rangs ; elle vint choisir entre nous tous celui que la Providence avait prédestiné pour être comme la première pierre de cet édifice céleste qui doit se former dans le séjour des élus sous le nom, ainsi que le P. Roger l'a si bien désigné, de Congrégation du ciel. Celui donc à qui il a été donné de devenir le noyau de cette réunion de nos amis qui ont enfin trouvé dans le sein de la charité le repos, la paix et la récompense et que nous espérons tous rejoindre un jour, celui qui le premier est venu attendre les autres fut M. Deroche, jeune professeur à l'école vétérinaire de Lyon. Il était d'une grande piété, sa vertu a été angélique ; il fut attaqué d'une fièvre muqueuse, dès son invasion la maladie se manifesta avec des caractères alarmants. Il était seul logé comme un garçon en chambre garnie. Il n'avait pour tous revenus, qu'un modeste appointement bien éloigné de pouvoir suffire aux soins que sa position exigeait. Dans cette occurrence la Congrégation dut tout naturellement le considérer comme sa famille ; elle se chargea de pourvoir à tout ce qui pouvait contribuer à le soulager et donna l'espérance de le rendre à la vie. Elle le fit transporter à l'Hôtel-Dieu dans une de ces
% Source : PDF original, vue 105, côté droit, page manuscrite 211.
chambres destinées aux malades en état de faire un peu de dépenses, et où, moyennant la somme de 12\up{f} par jour on peut réunir autour de soi les soins les plus minutieux et toutes les ressources de l'art aussi bien et mieux encore que dans la maison la plus opulente. Les congréganistes ne se bornèrent pas là ; ils prirent soin du malade en véritables frères ; jour et nuit ils se relevaient dans sa chambre pour lui rendre concurremment avec les frères et sœurs de l'hôpital tous les services qui dépendaient d'eux. Des soins si empressés et si touchants ne purent pas réussir à conserver dans ce monde ce bon ami ; son heure était venue, la récompense l'attendait. Nous vîmes donc avec douleur la maladie empirer de plus en plus. Il reçut en présence d'un grand nombre de congréganistes les sacrements de l'Église avec de grands sentiments de foi et de piété, et ce fut encore au milieu d'eux qu'il rendit son âme à son créateur et s'endormit dans le Seigneur.

La Congrégation se chargea de tous les frais des obsèques. Sans y mettre précisément du luxe elle voulut cependant les accompagner d'un peu d'appareil. Le corps fut présenté à l'église de l'hôpital, où on chanta une grand-messe, de là il fut transporté au cimetière de Loyasse. La Congrégation tout entière assistait à ses funérailles. Les congréganistes s'emparèrent du corps de leur frère et l'emportèrent sur leurs épaules pendant une partie du trajet, mais le convoi ayant été rencontré par l'école vétérinaire qui par suite d'une méprise se rendait à l'enterrement du professeur plus tard que l'heure indiquée, les élèves de cette école jaloux de témoigner leur affection au professeur qu'ils venaient de perdre enlevèrent à ses amis qui ne pouvaient déployer aucun caractère officiel
% Source : PDF original, vue 106, côté gauche, page manuscrite 212.
pour réclamer la préférence, le pieux fardeau dont ils s'étaient chargés et le portèrent à leur tour jusqu'au cimetière. Pour conserver le souvenir de notre premier né pour le ciel, la Congrégation fit placer une modeste pierre tumulaire sur la tombe de M. Deroche.

Une circonstance assez singulière accompagna la mort de ce pieux ami ; comme elle pourrait être considérée comme un indice de sa sainteté, il me paraît tout à fait convenable d'en rappeler le souvenir ; et cela d'autant mieux que ce que je vais vous raconter a été avéré dans le temps par le témoignage de plusieurs de ses amis qui ne permettent pas l'ombre d'un doute sur la vérité des faits. Quinze jours environ avant l'époque à laquelle M. Deroche tomba malade, il fit part dans l'intimité à un congréganiste avec lequel il était fort lié d'un rêve qu'il venait de faire et dont il était vivement frappé. Tout ce qui devait avoir lieu dans sa maladie, à sa mort et à son enterrement, avait été fidèlement représenté dans ce rêve ; tout absolument, même l'enlèvement de son corps des mains de ses amis que vint opérer l'école vétérinaire. Or pour deviner d'avance ce dernier détail il aurait fallu être averti que l'école vétérinaire ferait involontairement une erreur qui la mettrait en retard.

\section*{Mort de M. Terret l’aîné}
% Mention marginale des pages manuscrites 213–214 ; reportée au début du portrait.
Ce fut quelque temps après cette mort bienheureuse devant le Seigneur que nous perdîmes un autre congréganiste dont le nom se retrouve encore dans les listes de la Congrégation puisque son frère en fait partie. Sa mémoire ne s'est pas effacée non plus du cœur de ses anciens amis qui ne lisent pas ici sans
% Source : PDF original, vue 106, côté droit, page manuscrite 213.
éprouver quelque émotion le nom de Terret l'aîné. Doué d'un jugement droit, rempli de zèle pour la gloire de Dieu, gagnant tous les cœurs par l'aménité de ses manières, cet excellent ami avait toutes les qualités propres à attirer sur lui l'attention de la Congrégation. Aussi ce fut dès les premières années de son existence qu'elle prononça avec empressement son admission. Terret fut donc présenté en assemblée générale, mais si d'un côté il était animé du désir de faire le bien, de l'autre il n'aimait pas à s'engager à la légère, et de plus il possédait une fermeté d'âme qui n'est pas ordinaire. Lors donc qu'après lui avoir donné lecture du règlement, on lui demanda suivant l'usage s'il consentait à s'y soumettre, il répondit qu'il ne pouvait pas se déterminer ainsi ; qu'à la vérité on venait bien de lui donner connaissance de nos règles, mais qu'une lecture faite avec autant de rapidité ne pouvait pas lui suffire pour bien connaître les engagements qu'il devait prendre. Il conclut à ce qu'on voulût bien lui confier le règlement pour le méditer à son aise se réservant ensuite de rendre sa réponse. Cet incident qui ne s'était jamais présenté et qui à ce que je crois ne s'est pas renouvelé, embarrassa un peu le préfet. Il consulta les assistants et après en avoir délibéré avec eux il déclara au récipiendaire que nos usages s'opposaient absolument à ce qu'on pût ainsi remettre ce règlement entre des mains étrangères, mais que pour lui donner les éclaircissements qu'il désirait, tout ce qu'il était possible de faire en sa faveur était d'autoriser les assistants à lui en donner en particulier une seconde lecture. Terret fut parfaitement satisfait de ce moyen terme. Cette seconde lecture calma tous les scrupules et il accepta avec empressement.
% Source : PDF original, vue 107, côté gauche, page manuscrite 214.
La fidélité qu'il a mise jusqu'à la mort à remplir, non seulement avec exactitude, mais aussi avec zèle tous ses devoirs de congréganiste, a bien prouvé que cette espèce d'hésitation ne provenait absolument que de la pureté d'intention avec laquelle il agissait. Cette pureté d'intention lui faisait une loi de bien s'assurer de la nature des engagements qu'il allait prendre pour avoir une certitude complète de pouvoir y être fidèle. Pendant le temps malheureusement trop court que nous avons possédé cet excellent ami, il s'est distingué dans la Congrégation non seulement par sa bonne volonté mais encore par la capacité peu commune dont il était doué. Il a occupé les places de chef de section et d'assistant. Le préfet dans les circonstances difficiles se plaisait d'autant mieux à le consulter qu'il était plein de confiance dans ses lumières. Il s'est donc rendu de cette manière très utile à la Congrégation et de plus il était pour elle, par les vertus dont il lui donnait l'exemple un véritable et un grand sujet d'édification. Après l'avoir ainsi édifiée pendant sa vie, il l'a peut-être édifiée davantage encore par la patience avec laquelle il a supporté sa longue maladie et s'est résigné à la sainte mort qui lui a ouvert encore bien jeune les portes de l'éternité.

Le souvenir de la perte de ces amis dont nos cœurs ont été successivement affligés et la pensée plus consolante des vertus dont ils nous ont rendu les témoins, nous ont éloignés pour quelques instants des grands événements qui en troublant la paix de l'Église faisaient en même temps éclore dans notre Congrégation des inquiétudes trop bien fondées. Il est temps d'y revenir car ils méritent toute notre attention.

% Source : PDF original, vue 107, côté droit, page manuscrite 215.
Les passions brûlantes qui avaient excité l'ambitieux Napoléon à persécuter le chef de l'Église s'envenimaient de plus en plus. Le pape avait formellement refusé son consentement au divorce de l'empereur. Celui-ci s'était fait donner une consultation par quelques théologiens complaisants, et s’en était contenté. L'empereur d'Allemagne aveuglé par cette fumée de gloire qui environne Napoléon et craignant aussi de descendre du trône accepte l'alliance humiliante qui lui est offerte ; il accorde au soldat parvenu la main de la fille des Césars. Ceux des cardinaux qui fidèles à leurs devoirs, refusent d'assister à ce mariage illicite, sont par l'ordre du despote dépouillés de la pourpre romaine ; distinction honorable qui leur valut le glorieux surnom de cardinaux noirs ; ce qui peut littéralement se traduire par ces mots : cardinaux fidèles.

\section*{Pie VII enlevé de Rome}
% Mention marginale de la page manuscrite 215.
Napoléon ne s'en tiendra pas là ; il faut que son orgueil humilié se venge par des violences ; les troupes impériales reçoivent l'ordre de marcher en avant et de s'emparer de la ville de Rome, on envahit le palais du souverain pontife ; on enlève le chef de l'Église ; comme son prédécesseur Pie VI il est traîné brutalement en captivité. L'empereur le fait conduire à Savonne ; là il le sépare des cardinaux qui l'avaient accompagné et les fait jeter dans les fers.

\section*{Indulgences}
% Mention marginale de la page manuscrite 216 ; reportée au début du récit.
La veille du jour fatal où il fut donné à l'ennemi de l'Église d'éprouver de nouveau par la persécution le vicaire de J. C., ce père commun des fidèles, en promenant par la pensée ses regards paternels sur la multitude d'Églises du monde catholique dont le soin lui a été confié, jette un coup d'œil sur cette terre des martyrs.
% Source : PDF original, vue 108, côté gauche, page manuscrite 216.
dans laquelle l'arbre de la foi arrosé par le sang des Pothin, des Irénées et de ces invincibles légions qui à leur suite ont triomphé par une mort glorieuse de toute la rage de l'enfer, a poussé de profondes racines ; il aperçoit le point presque imperceptible où se trouve cette Congrégation lyonnaise, toujours dévouée à la chaire de St Pierre ; il la considère gémissant sur les malheurs de la religion et suppliant le Seigneur de mettre à cette nouvelle persécution un terme prochain. Aussitôt ses entrailles paternelles sont émues, il ouvre les trésors de l'Église, il les répand non seulement avec abondance mais presque avec profusion sur ses enfants chéris.

Oui, Messieurs, reconnaissance éternelle à ce vénérable pontife Pie VII. C'est au milieu de toutes les douleurs, de toutes les angoisses dont son cœur est abreuvé que la veille de son enlèvement il nous a fait expédier de Rome ce bref d'indulgences qui nous en a si largement pourvus. Non seulement depuis lors les congréganistes ont eu la faculté de gagner de fréquentes indulgences plénières, mais dans le cours habituel de la vie, la plus légère bonne œuvre opérée par eux a trouvé sa récompense dans une indulgence de sept ans et de sept quarantaines. Il faut observer ici que dans la plupart des brefs d'indulgences celles accordées ainsi pour la pratique d'une bonne œuvre, n'excèdent guère ordinairement soixante à cent jours, au plus trois cents jours. Dans le bref dont nous sommes maintenant en possession et dont nous aurons à parler plus tard, toutes ces faveurs spirituelles nous ont été fidèlement conservées.

C'est ainsi que le chef visible de toutes les Églises, semblable à un bon père de famille qui dans quelque position qu'il
% Source : PDF original, vue 108, côté droit, page manuscrite 217.
se trouve, le souvenir de ses enfants chéris n'abandonne jamais, s'occupait dans ces moments critiques pour nous devant le Seigneur, et usant dans sa plénitude de ce pouvoir de délier que J. C. a remis entre ses mains, il s'efforçait de nous faciliter les moyens d'expier nos fautes et adoucissait en notre faveur la pente rapide qu'il faut nécessairement gravir pour monter au ciel.

Ce fut pour nous une bien grande consolation que cette nouvelle marque de bonté par laquelle le Saint-Père encourageait nos efforts ; jamais il n'avait été plus nécessaire de s'armer de courage ; car l'horizon se rembrunissait de plus en plus et la persécution déjà commencée tout autour de nous, faisant supporter ses rigueurs au chef de l'Église, ne pouvait tarder à nous atteindre.

Au milieu de ces circonstances si alarmantes soit pour l'Église en général soit en particulier pour notre Congrégation, arriva l'expiration de la troisième année d'exercice de M. Philpin en qualité de préfet. Il fallut donc, malgré le désir que la Congrégation aurait eu de le conserver plus longtemps à sa tête, se conformer au règlement et lui donner un successeur. M. Franchet fut appelé dans ces circonstances difficiles à gouverner la Congrégation.
